% Appendix C -- KITTI supplementary results (revision v2: KITTI part of the original Appendix A + KITTI setup details from the main text,
% best.pt artifact analysis (former 8.6(b)), H07 KITTI details and latency bug (former 8.6(c))).
% Included via \input from esm.tex.
\setcounter{table}{0}\setcounter{figure}{0}
\section{KITTI supplementary results}\label{sec:ek-ikincil}

This section gives the setup details of the KITTI experiment (Section~\ref{sec:genelleme} of the main article) and its secondary results. It also covers the training curve of H05, the analysis of \texttt{best.pt} selection, the post-unlock robustness analysis, and the details of training-free inference.

\subsection{Setup details}\label{sec:ek-kitti-kurulum}

\paragraph{Questions, rules, and training details.}
S1 is answered per model with two families. The improvement family counts the cards that meet the acceptance criterion. The directional replication family asks whether the loss seen on VOC is replicated. S2 is answered with the contribution and additivity criteria of the two component cards run on VOC. Later changes to the rule document do not apply to the VOC cards of S2. The actual optimizer behavior on VOC was fixed with the following settings: MuSGD, learning rate 0.01, momentum 0.937, bias warmup 0. The copy-paste parameters were unchanged. The number of objects pasted per mosaic is 0.98 in H03 and 2.00 in H08. In YOLO26n, the H02 and H05 patches are applied to both training heads; a head counter verified this. With three seeds, the critical effect at 50\% power in the first step of Holm-5 is $\approx2.8\,s_\Delta$. For a card that is not accepted, we state ``acceptance criterion not met''. In H07, the NMS, Soft-NMS, and WBF rows are not applicable to YOLO26n, which is NMS-free.

\paragraph{Split and ignore rule.}
We did not use the ready-made Ultralytics split. Nor did we use the 3712/3769 train/validation split common in 3D detection work. That split has no separate test partition. We also could not verify that it includes a check (NCC) that groups frames of the same scene stored under different recording names. Scene groups were assigned to the training, validation, and test partitions under two constraints. At least 40\%, 5\%, and 10\% of each class's objects fall in the respective partition. No single group exceeds 25\% of the target validation or test size. The highest cross-partition NCC is 0.855 for test--train and 0.888 for validation--train. The test partition contains 6\,212 objects and 2\,436 \texttt{DontCare} regions. In evaluation, \texttt{DontCare} regions are added for every class as COCO-style ignore regions (\texttt{iscrowd}). Detections are first matched to real objects. Only an unmatched detection can be ignored. It is ignored if its overlap with a \texttt{DontCare} region is greater than or equal to the current IoU threshold. Overlap is measured as IoA, the intersection divided by the detection area. The official KITTI AP, by contrast, uses a fixed per-class threshold and difficulty levels. mAP without \texttt{DontCare} is reported as a secondary result in Table~\ref{tab:kitti-ikincil}.

\paragraph{Training, health gate, and learning check.}
The loss code of YOLO26n is v8.4.165 \texttt{E2ELoss}. The number of iterations on KITTI is $\approx$3.1 times lower than on VOC (24\,300). The head counter was read in every H02 and H05 run. A training run took 1.68--1.83 GPU hours for YOLO11s and 1.17--1.31 for YOLO26n. Every run passes an automatic gate with three parts. The batch must not be silently reduced, the patch must actually be called, and a learning check must pass. The learning check is a validity criterion with five conditions:
\texttt{results.csv} contains exactly 100 epochs;
loss and metric values are finite;
the learning rate schedule matches that of the seed-0 baseline of the same model, epoch by epoch;
the mean training loss (box $+$ classification) over epochs 81--90 is lower than over epochs 11--20;
and the best validation mAP50-95 is above 0.02.
No threshold relative to the baseline was set, because it would censor strongly negative cards. A run that fails does not enter $\Delta$. A best epoch below 20 is a warning, not an invalidation. The rule is applied identically to all cards and both models. Two rule changes were made: this check was added to the rule document, and \texttt{last.pt} evaluation was added to the test script (Section~\ref{sec:sinirliliklar} of the main article). Until the YOLO11s queue finished, monitoring used only the event type, run ID, and exit code.

\paragraph{Resolution.}
Most KITTI frames are 1242$\times$375 (1\,396/1\,504 in the test set), with an aspect ratio of 3.31. At 640 px input, the long side is scaled by $\approx$0.52, so a frame becomes $\approx$640$\times$193. In the test set, 19.1\% of pedestrian and cyclist objects fall in the COCO ``small'' category (area $<32^2$ px) at native resolution. At 640 px, this share is 65.6\%. VOC images, in contrast, are at most 500 px and are upscaled by a factor of 1.28 at 640 px. As a result, the small-object share in the VOC test set drops from 4.7\% to 1.85\%. Following the single-variable rule, the training resolution was kept the same as for VOC (640 px). This difference is the simplest hypothesis for two contrasts. First, 832 px test input yields a gain on KITTI but a loss on VOC. Second, multi-scale training (H01) is accepted on KITTI but not on VOC. The reasoning is that KITTI has a high small-object share, so a higher effective resolution may help. On VOC, the objects are already upscaled at 640 px. The hypothesis was not tested.

\paragraph{Baseline plausibility.}
No published anchor exists for KITTI with this split and evaluation rule. The VOC plausibility gate was therefore not applied, and three checks were made instead. The first is a training pipeline diagnosis. The baseline \texttt{last.pt} weights give a high mAP50-95 on training images and a low one on validation (numbers in Section~\ref{sec:ek-kontrol-noktasi}). The second check is the absolute values, in points, with the \texttt{DontCare} rule and averaged over three seeds. Test mAP50-95 is 19.30 (\texttt{best.pt}) and 19.66 (\texttt{last.pt}) for YOLO11s, and 18.85 and 17.35 for YOLO26n. The \texttt{best.pt} validation mAP50-95 is 19.26 $\pm$ 2.10 for YOLO11s and 18.51 $\pm$ 1.49 for YOLO26n. Test car AP50-95 (\texttt{best.pt}) is 57.66 for YOLO11s and 58.10 for YOLO26n. These values only indicate the scale of the $\Delta$s. They are not a comparison anchor, and they are not compared with the VOC values or across models. The third check is an external reference value, used for order of magnitude only (Section~\ref{sec:ek-dis}). Its resolution, split, and evaluation differ from ours, so it is not comparable with our values.

\subsection{KITTI secondary results and best epoch per run}
This subsection gives the secondary reports defined in the KITTI rule document. All are descriptive. There is no Holm correction and no significance claim.

% generator: evidence/scripts/build_ek_tables.py <- evidence/evidence.json (ek.test.*, ek.val.*, ek3.test.*, ek3.b2.*; manual edits are lost at the next regeneration)
\begin{table}[!htbp]
\centering\small\setlength{\tabcolsep}{6pt}
\caption{KITTI pre-specified secondary descriptive comparisons: card $-$ baseline $\Delta$, mean $\pm$ s.d. [points], $n=3$ seed-paired. No Holm correction, no color.}\label{tab:kitti-ikincil}
\begin{tabularx}{\textwidth}{>{\raggedright\arraybackslash}X c c c c}
\toprule
\multicolumn{1}{@{}l}{} & \multicolumn{3}{c}{Test $\Delta$} & Val $\Delta$ \\
\cmidrule(lr){2-4}
Card & No DontCare & 3 classes & 7 classes & ep11--90 curve \\
\midrule
\multicolumn{5}{@{}l}{\textbf{YOLO26n}} \\
H01 multi-scale & +0.05 $\pm$ 2.19 & +1.22 $\pm$ 1.18 & +0.08 $\pm$ 2.52 & $-$0.31 $\pm$ 0.03 \\
H02 Inner-CIoU & $-$1.01 $\pm$ 1.98 & +1.04 $\pm$ 1.17 & $-$1.02 $\pm$ 2.29 & +0.51 $\pm$ 0.52 \\
H03 box copy-paste & +0.88 $\pm$ 0.68 & +2.81 $\pm$ 0.65 & +1.07 $\pm$ 0.87 & +0.62 $\pm$ 1.02 \\
H04 mosaic scale & $-$1.33 $\pm$ 2.19 & $-$0.97 $\pm$ 1.38 & $-$1.53 $\pm$ 2.46 & $-$0.54 $\pm$ 0.48 \\
H05 Varifocal & +0.47 $\pm$ 1.48 & +1.28 $\pm$ 1.16 & +0.78 $\pm$ 1.73 & +1.29 $\pm$ 0.20 \\
H06 box loss gain & $-$0.24 $\pm$ 0.72 & +0.56 $\pm$ 1.68 & $-$0.26 $\pm$ 0.86 & +0.33 $\pm$ 0.71 \\
H08 high dose & +0.67 $\pm$ 0.68 & +2.68 $\pm$ 1.32 & +0.73 $\pm$ 0.89 & +1.01 $\pm$ 0.83 \\
\midrule
\multicolumn{5}{@{}l}{\textbf{YOLO11s}} \\
H01 multi-scale & +1.48 $\pm$ 0.39 & +0.48 $\pm$ 1.72 & +1.72 $\pm$ 0.29 & +0.59 $\pm$ 0.57 \\
H02 Inner-CIoU & +1.13 $\pm$ 0.48 & +0.06 $\pm$ 1.75 & +1.38 $\pm$ 0.66 & $-$0.27 $\pm$ 0.55 \\
H03 box copy-paste & +1.17 $\pm$ 1.00 & +0.22 $\pm$ 2.73 & +1.50 $\pm$ 1.17 & +1.01 $\pm$ 0.55 \\
H04 mosaic scale & +1.00 $\pm$ 0.83 & $-$0.66 $\pm$ 1.86 & +0.86 $\pm$ 1.30 & $-$0.36 $\pm$ 1.36 \\
H05 Varifocal & +2.35 $\pm$ 0.96 & +1.41 $\pm$ 2.06 & +2.73 $\pm$ 1.60 & +0.53 $\pm$ 0.37 \\
H06 box loss gain & +0.45 $\pm$ 0.65 & +0.13 $\pm$ 1.31 & +0.57 $\pm$ 0.79 & $-$0.38 $\pm$ 0.29 \\
H08 high dose & +1.63 $\pm$ 2.89 & +1.30 $\pm$ 4.20 & +1.66 $\pm$ 4.04 & +1.00 $\pm$ 0.39 \\
\bottomrule
\addlinespace[3pt]
\multicolumn{5}{@{}p{\dimexpr\textwidth-2\tabcolsep\relax}@{}}{Test (\texttt{best.pt}): mAP50-95 without DontCare. Subsets with the DontCare rule: car, pedestrian, cyclist (3 classes) and 7 classes excluding person\_sitting. Val: mean of the training curve (Ultralytics, without DontCare) over epochs 11--90. Test object counts: person\_sitting 56, tram 126, truck 119.} \\
\end{tabularx}
\end{table}


Table~\ref{tab:kitti-ikincil} shows the test $\Delta$ without \texttt{DontCare}. In none of the 14 cards does it differ from the primary $\Delta$ (with the \texttt{DontCare} rule) by more than 0.24 points. The primary result is therefore not sensitive to the ignore rule. On the standard three classes (car, pedestrian, cyclist), the gain of YOLO11s H01 is $+0.48\pm1.72$. Its per-seed $\Delta$s are +1.57, $-$1.50, and +1.38. On the seven-class subset that excludes person\_sitting, the gain is $+1.72\pm0.29$, positive in all three seeds. The difference appears when van, truck, tram, and misc are added to the standard three classes. These classes have far fewer test objects than car (4\,457): van has 464, misc 184, tram 126, and truck 119. person\_sitting has 56. With these counts, per-class AP is noisy. The distribution of the gain over individual classes cannot be resolved with these data. YOLO26n gives an example in the opposite direction. Its three-class difference is $+2.81\pm0.65$ for H03 (positive in all three seeds) and $+2.68\pm1.32$ for H08. The primary $\Delta$s of these two cards are far below this (Table~\ref{tab:kitti-test} of the main article). The validation curve column uses the mean validation mAP50-95 over epochs 11--90 during training. It reports the difference of this mean from the baseline. For H05, it is $+1.29\pm0.20$ for YOLO26n and $+0.53\pm0.37$ for YOLO11s (Section~\ref{sec:ek-h05-kitti}).

% generator: evidence/scripts/build_ek_tables.py <- evidence/evidence.json (ek.epoch.*, ek3.val.*; manual edits are lost at the next regeneration)
\begin{table}[!htbp]
\centering\small\setlength{\tabcolsep}{4pt}
\caption{KITTI: best epoch per run (\texttt{best.pt}) and val best$-$last difference [points]. 48 runs, seeds s0--s2. Italic: epoch 1.}\label{tab:kitti-epoch}
\begin{tabularx}{\textwidth}{>{\raggedright\arraybackslash}X c c c c c c}
\toprule
\multicolumn{1}{@{}l}{} & \multicolumn{2}{c}{s0} & \multicolumn{2}{c}{s1} & \multicolumn{2}{c}{s2} \\
\cmidrule(lr){2-3}\cmidrule(lr){4-5}\cmidrule(lr){6-7}
Arm & Epoch & Diff. & Epoch & Diff. & Epoch & Diff. \\
\midrule
\multicolumn{7}{@{}l}{\textbf{YOLO26n}} \\
B baseline & 16 & +4.3 & 11 & +5.1 & 28 & +3.7 \\
H01 multi-scale & 24 & +3.8 & 9 & +4.6 & 20 & +5.3 \\
H02 Inner-CIoU & 29 & +3.8 & 11 & +4.0 & 38 & +3.1 \\
H03 box copy-paste & 10 & +4.7 & 24 & +4.4 & 37 & +2.6 \\
H04 mosaic scale & 16 & +5.5 & 11 & +4.0 & 12 & +5.4 \\
H05 Varifocal & 46 & +2.9 & 39 & +4.2 & 33 & +2.8 \\
H06 box loss gain & 51 & +2.9 & 11 & +4.8 & 28 & +3.4 \\
H08 high dose & 39 & +3.4 & 12 & +4.6 & 48 & +3.3 \\
\midrule
\multicolumn{7}{@{}l}{\textbf{YOLO11s}} \\
B baseline & 9 & +4.4 & 32 & +3.1 & \textit{1} & +5.7 \\
H01 multi-scale & 15 & +4.6 & 48 & +3.9 & \textit{1} & +5.5 \\
H02 Inner-CIoU & 29 & +4.6 & 38 & +4.6 & \textit{1} & +6.6 \\
H03 box copy-paste & \textit{1} & +4.8 & 12 & +2.7 & \textit{1} & +5.2 \\
H04 mosaic scale & 39 & +5.1 & \textit{1} & +7.4 & \textit{1} & +6.1 \\
H05 Varifocal & 47 & +2.7 & 54 & +1.9 & 17 & +3.2 \\
H06 box loss gain & 13 & +4.9 & \textit{1} & +3.6 & \textit{1} & +6.5 \\
H08 high dose & \textit{1} & +2.9 & \textit{1} & +6.3 & 24 & +3.9 \\
\bottomrule
\addlinespace[3pt]
\multicolumn{7}{@{}p{\dimexpr\textwidth-2\tabcolsep\relax}@{}}{Val difference: in-training best $-$ last-epoch mAP50-95 (Ultralytics, without DontCare). Runs whose best epoch is 1: YOLO26n 0/24, YOLO11s 11/24.} \\
\end{tabularx}
\end{table}


Table~\ref{tab:kitti-epoch} gives the best epoch of each run. It also gives the in-training validation difference between the best and the last epoch. For YOLO11s, the best epoch is 1 in 11 of 24 runs: 2 in seed 0, 3 in seed 1, and 6 in seed 2. For YOLO26n, this happens in no run. The difference is positive in all 48 runs. This is expected, because \texttt{best.pt} is selected on validation. The difference is measured with Ultralytics' in-training evaluator, which does not use \texttt{DontCare}. The 24-run mean is $+4.03$ for YOLO26n and $+4.58$ for YOLO11s. For the baseline alone, it is $+4.37$ and $+4.39$ points. These numbers come from a different measurement than the $+3.23$ and $+3.84$ in Section~\ref{sec:ek-kontrol-noktasi}, which come from this study's COCO-style re-evaluation.

% generator: evidence/scripts/build_ek_tables.py <- evidence/evidence.json (ek4.test.*, ek4.val.*, ek.val.*; manual edits are lost at the next regeneration)
\begin{table}[!htbp]
\centering\small\setlength{\tabcolsep}{3pt}
\caption{KITTI secondary metrics (\texttt{best.pt}): card $-$ baseline $\Delta$, mean $\pm$ s.d. [points], $n=3$ seed-paired. Test uses the DontCare rule. Descriptive, no color.}\label{tab:kitti-test-ek}
\begin{tabularx}{\textwidth}{>{\raggedright\arraybackslash}X c c c c c c}
\toprule
\multicolumn{1}{@{}l}{} & \multicolumn{4}{c}{Test $\Delta$} & \multicolumn{2}{c}{Val [points]} \\
\cmidrule(lr){2-5}\cmidrule(lr){6-7}
Card & mAP50 & AP$_S$ & AP$_M$ & AP$_L$ & $\Delta$ mAP50-95 & In-train$-$re-eval. \\
\midrule
\multicolumn{7}{@{}l}{\textbf{YOLO26n}} \\
B & 33.59$\pm$0.89 & 7.92$\pm$2.61 & 15.94$\pm$1.05 & 32.71$\pm$0.94 & -- & +0.35$\pm$1.20 \\
H01 & +0.23$\pm$1.82 & +0.05$\pm$1.21 & $-$0.37$\pm$1.48 & +2.32$\pm$3.38 & $-$0.10$\pm$2.25 & +0.21$\pm$1.01 \\
H02 & $-$1.31$\pm$2.54 & +0.29$\pm$2.68 & $-$0.51$\pm$2.18 & $-$1.86$\pm$0.40 & $-$0.29$\pm$0.80 & +0.51$\pm$0.52 \\
H03 & +1.94$\pm$0.43 & $-$0.10$\pm$1.47 & +1.05$\pm$0.18 & +1.81$\pm$3.11 & +0.63$\pm$1.06 & +0.19$\pm$0.46 \\
H04 & $-$1.90$\pm$2.14 & $-$1.80$\pm$1.65 & $-$0.83$\pm$2.08 & $-$1.38$\pm$0.97 & $-$0.22$\pm$1.89 & +0.78$\pm$0.59 \\
H05 & +0.40$\pm$2.14 & +0.13$\pm$1.82 & +1.02$\pm$1.75 & +0.57$\pm$2.08 & +0.68$\pm$1.14 & +0.44$\pm$0.53 \\
H06 & $-$0.38$\pm$1.10 & $-$0.30$\pm$1.92 & +0.18$\pm$1.31 & $-$0.90$\pm$1.91 & +0.09$\pm$0.65 & +0.20$\pm$0.32 \\
H08 & +1.18$\pm$0.69 & $-$1.24$\pm$3.08 & +0.72$\pm$1.54 & +0.96$\pm$1.53 & +0.46$\pm$1.34 & +0.29$\pm$0.52 \\
\midrule
\multicolumn{7}{@{}l}{\textbf{YOLO11s}} \\
B & 35.39$\pm$0.65 & 7.20$\pm$0.91 & 18.14$\pm$0.80 & 31.92$\pm$0.86 & -- & +1.12$\pm$0.68 \\
H01 & +1.63$\pm$1.08 & +2.58$\pm$1.61 & +0.98$\pm$1.77 & +2.74$\pm$2.76 & +1.73$\pm$1.78 & $-$0.01$\pm$0.33 \\
H02 & +1.87$\pm$2.46 & +1.28$\pm$0.84 & +0.95$\pm$0.84 & +2.10$\pm$1.24 & +1.15$\pm$1.05 & +0.28$\pm$0.29 \\
H03 & +1.84$\pm$2.56 & +2.50$\pm$2.35 & +0.60$\pm$1.22 & +1.28$\pm$1.94 & +1.86$\pm$2.44 & +0.65$\pm$0.54 \\
H04 & +2.14$\pm$0.26 & +2.81$\pm$3.12 & +0.81$\pm$1.34 & +2.50$\pm$1.03 & +1.82$\pm$0.43 & +0.74$\pm$1.00 \\
H05 & +1.66$\pm$0.76 & +4.37$\pm$2.90 & +1.60$\pm$1.01 & +2.30$\pm$1.57 & +0.33$\pm$2.09 & +0.47$\pm$0.33 \\
H06 & +0.80$\pm$1.62 & +5.69$\pm$1.31 & +0.62$\pm$0.40 & +2.10$\pm$1.00 & +1.24$\pm$0.56 & +0.03$\pm$0.38 \\
H08 & +2.04$\pm$1.65 & +1.99$\pm$1.28 & +0.58$\pm$0.85 & +4.31$\pm$1.98 & +2.21$\pm$3.24 & $-$0.06$\pm$0.61 \\
\bottomrule
\addlinespace[3pt]
\multicolumn{7}{@{}p{\dimexpr\textwidth-2\tabcolsep\relax}@{}}{Cards as in Table~\ref{tab:kitti-test}. Val $\Delta$: mAP50-95 with the DontCare rule. In-train$-$re-eval.: in-training best val (Ultralytics) $-$ COCO-style re-evaluation of the same weights without DontCare. Row B (baseline) gives absolute values. AP$_S$/AP$_M$/AP$_L$ are the COCO area classes.} \\
\end{tabularx}
\end{table}


Table~\ref{tab:kitti-test-ek} gives three quantities per card. These are the secondary test metrics, the validation $\Delta$, and the difference between the in-training evaluation and the re-evaluation. Row B gives absolute values. For YOLO11s, seven of seven cards are positive on all four metrics. The ranges are $+0.80$ to $+2.14$ for mAP50, $+1.28$ to $+5.69$ for AP$_S$, $+0.58$ to $+1.60$ for AP$_M$, and $+1.28$ to $+4.31$ for AP$_L$. The validation $\Delta$ is also positive for all seven cards. This fits a common difference that shifts all cards against the same baseline better than a card-specific effect (Section~\ref{sec:ek-kontrol-noktasi}). For YOLO26n, the signs are mixed. The number of positive cards is 4/7 for mAP50, 3/7 for AP$_S$, 4/7 for AP$_M$, and 4/7 for AP$_L$. The last column compares two evaluations of the same \texttt{best.pt} weights. It is the best in-training validation value minus the COCO-style re-evaluation without \texttt{DontCare}. The 24-run mean is $+0.37$ for YOLO26n and $+0.40$ for YOLO11s. Over all 48 runs, the range is $-1.03$ to $+1.80$, and 37/48 runs are positive. \texttt{best.pt} is selected by Ultralytics fitness. It therefore relies on a different criterion from the evaluator of the primary endpoint. AP$_S$ is the COCO small-object category. For YOLO26n, the s.d.s in this column range from 1.21 to 3.08. Individual card differences should be read with this noise in mind.

% generator: evidence/scripts/build_ek_tables.py <- evidence/evidence.json (ek4.test.*.sinif.*, ek3.b2.test_nesne.*; manual edits are lost at the next regeneration)
\begin{table}[!htbp]
\centering\small\setlength{\tabcolsep}{2pt}
\caption{KITTI test: per-class AP50-95 (DontCare rule, \texttt{best.pt}), card $-$ baseline $\Delta$, mean $\pm$ s.d. [points], $n=3$ seed-paired. Descriptive, no color.}\label{tab:kitti-sinif}
\begin{tabularx}{\textwidth}{>{\raggedright\arraybackslash}X c c c c c c c c}
\toprule
Card & car & van & truck & ped. & sit. & cyc. & tram & misc \\
\midrule
\multicolumn{9}{@{}l}{\textbf{YOLO26n}} \\
B & 58.1 & 12.4 & 13.7 & 24.5 & 1.1 & 7.2 & 33.3 & 0.6 \\
H01 & $-$0.6$\pm$2.0 & +0.2$\pm$1.9 & $-$1.6$\pm$4.0 & +1.8$\pm$2.0 & +0.5$\pm$0.5 & +2.5$\pm$2.6 & $-$1.5$\pm$11.0 & $-$0.1$\pm$0.4 \\
H02 & +0.3$\pm$1.2 & $-$1.6$\pm$1.9 & $-$3.3$\pm$3.0 & +0.4$\pm$2.4 & $-$0.3$\pm$0.8 & +2.4$\pm$0.8 & $-$5.5$\pm$12.3 & +0.1$\pm$0.3 \\
H03 & +0.5$\pm$1.3 & +1.8$\pm$0.4 & +1.4$\pm$2.9 & +1.8$\pm$0.6 & +0.6$\pm$0.9 & +6.1$\pm$1.8 & $-$4.3$\pm$5.2 & +0.3$\pm$0.3 \\
H04 & $-$1.3$\pm$1.3 & $-$0.5$\pm$2.5 & $-$0.3$\pm$1.7 & $-$2.4$\pm$2.0 & +0.8$\pm$1.1 & +0.7$\pm$1.3 & $-$7.2$\pm$10.7 & +0.2$\pm$0.5 \\
H05 & +1.3$\pm$0.4 & +0.8$\pm$1.8 & $-$3.2$\pm$4.3 & +0.9$\pm$0.5 & $-$0.5$\pm$0.1 & +1.6$\pm$2.8 & +3.1$\pm$5.9 & +0.8$\pm$0.7 \\
H06 & $-$0.3$\pm$1.1 & $-$0.6$\pm$1.3 & $-$0.4$\pm$5.6 & $-$0.1$\pm$1.2 & $-$0.2$\pm$0.8 & +2.1$\pm$3.6 & $-$2.6$\pm$8.6 & +0.1$\pm$0.3 \\
H08 & +0.7$\pm$1.6 & $-$0.8$\pm$2.0 & +1.0$\pm$4.4 & +2.2$\pm$1.4 & +0.6$\pm$0.7 & +5.2$\pm$1.1 & $-$3.0$\pm$5.9 & $-$0.1$\pm$0.4 \\
\midrule
\multicolumn{9}{@{}l}{\textbf{YOLO11s}} \\
B & 57.7 & 12.8 & 13.5 & 27.0 & 3.4 & 15.7 & 22.9 & 1.4 \\
H01 & +2.2$\pm$0.9 & +2.4$\pm$0.9 & +0.1$\pm$3.0 & +1.8$\pm$2.4 & $-$0.1$\pm$0.7 & $-$2.6$\pm$3.4 & +8.2$\pm$3.0 & $-$0.2$\pm$0.5 \\
H02 & +1.0$\pm$1.0 & $-$0.4$\pm$3.0 & +1.6$\pm$3.2 & +1.1$\pm$0.7 & $-$0.6$\pm$0.7 & $-$1.9$\pm$4.9 & +7.8$\pm$5.6 & +0.5$\pm$0.6 \\
H03 & +1.7$\pm$1.6 & +2.8$\pm$0.5 & +3.9$\pm$3.4 & +1.9$\pm$4.2 & +0.7$\pm$5.1 & $-$3.0$\pm$3.2 & +2.9$\pm$4.1 & +0.3$\pm$1.6 \\
H04 & $-$0.6$\pm$2.6 & +3.1$\pm$5.2 & +1.1$\pm$4.5 & $-$0.3$\pm$1.0 & +2.5$\pm$4.8 & $-$1.1$\pm$4.6 & +3.5$\pm$11.8 & +0.3$\pm$1.1 \\
H05 & +4.3$\pm$1.8 & +2.2$\pm$2.5 & +4.1$\pm$2.5 & +1.7$\pm$0.5 & $-$0.8$\pm$5.4 & $-$1.8$\pm$4.3 & +8.3$\pm$6.8 & +0.2$\pm$1.1 \\
H06 & +0.0$\pm$2.1 & +1.5$\pm$3.9 & +3.3$\pm$5.1 & +2.6$\pm$2.9 & +0.6$\pm$1.2 & $-$2.1$\pm$4.9 & $-$1.1$\pm$10.9 & $-$0.1$\pm$1.0 \\
H08 & +2.1$\pm$4.0 & +2.6$\pm$1.9 & +1.4$\pm$4.3 & +4.0$\pm$2.4 & +1.9$\pm$8.5 & $-$2.2$\pm$6.3 & +2.7$\pm$21.6 & +0.9$\pm$1.3 \\
\bottomrule
\addlinespace[3pt]
\multicolumn{9}{@{}p{\dimexpr\textwidth-2\tabcolsep\relax}@{}}{Cards as in Table~\ref{tab:kitti-test}. Abbreviation: sit. = person\_sitting. Row B gives absolute AP. Test object counts: car 4457, van 464, truck 119, ped. 574, sit. 56, cyc. 232, tram 126, misc 184.} \\
\end{tabularx}
\end{table}


Table~\ref{tab:kitti-sinif} gives the per-class AP50-95 $\Delta$. The other classes have far fewer test objects than car; the counts are given above and in the table footnote. For classes with few samples, the s.d.s can therefore exceed the differences themselves. For example, truck in YOLO11s H01 gives $+0.1\pm3.0$. The distribution of the gain over classes cannot be resolved with these data.

\subsection{Training curve of H05 on KITTI}\label{sec:ek-h05-kitti}
On KITTI, the same intervention did not produce the loss seen on VOC (Section~\ref{sec:genelleme-sonuc} of the main article). Measured: H05 peaks later than the baseline in both models. For YOLO11s, the best epochs are 47/54/17 for H05 and 9/32/1 for the baseline. For YOLO26n, they are 46/39/33 for H05 and 16/11/28 for the baseline. The peak is later in 6/6 seeds. This instance is weak for YOLO11s seed 2, however, because the baseline's best epoch there is 1. The late peak is most consistent for H05. For each of the other six cards, a later peak than the baseline appears in at most 4 of the six seeds of the two models (H02, H08). For YOLO26n, H05 has the largest difference among the cards in the epoch 11--90 mean of the validation curve. This difference is +1.29 $\pm$ 0.20 and is positive in all three seeds. For YOLO11s, it is not specific to H05. H05 gives +0.53 $\pm$ 0.37, while H03 gives +1.01 and H08 gives +1.00 (Table~\ref{tab:kitti-ikincil}). The mechanism was active in every run. For the remaining hypothesis and its limitation, see Section~\ref{sec:genelleme-sonuc} of the main article (answer to S2).

\subsection{Checkpoint selection: \texttt{best.pt} analysis}\label{sec:ek-kontrol-noktasi}
This subsection details the sensitivity to \texttt{best.pt} selection that is summarized in Section~\ref{sec:genelleme-sonuc} of the main article. The analysis is exploratory. The symptom is that the YOLO11s \texttt{best.pt} $\Delta$ is positive in seven of seven cards (Fig.~\ref{fig:best-last}b). This includes H02 and H05, which are downward on VOC. Four candidate explanations were tested, with the outcome given below:
an early peak;
noise in the validation set of only 10 scene groups;
the outlying baseline seed 2;
and the selection of \texttt{best.pt} by validation fitness without \texttt{DontCare}, while the primary endpoint uses the \texttt{DontCare} rule.

{\pgfplotsset{every axis post/.append style={height=4.2cm}}% display height only
% F7 -- KITTI test: per-card best.pt and last.pt Δ mAP50-95 (3-seed mean ± 1 s.d.), one panel per model (stacked).
% Data: figures/data/f7a_yolo26n.dat, f7b_yolo11s.dat (generator: figures/scripts/f7_best_last.py <- evidence/evidence.json ek.test.*).
% Non-color channel: best.pt = filled circle (Holm-significant improvement = filled triangle, blue), last.pt = open square.
\begin{figure}[!htbp]
\centering
\pgfplotstableset{col sep=semicolon}
\pgfplotsset{
  fseven/.style={scale only axis, width=0.52\textwidth, xmin=-4.6, xmax=5.0, ymin=0.4, ymax=7.6, ytick=data, font=\small,
    tick label style={font=\small}, label style={font=\small}, title style={font=\small, yshift=-1pt},
    xmajorgrids=true, grid style={black!10}, axis line style={black!70}, clip=false, every axis plot/.append style={forget plot}},
  fseven best/.style={only marks, mark=*, mark size=2.4pt, draw=black, fill=black,
    /pgfplots/error bars/x dir=both, /pgfplots/error bars/x explicit, /pgfplots/error bars/error bar style={line width=1.4pt, black}, /pgfplots/error bars/error mark options={rotate=90, mark size=2.6pt, black}},
  fseven good/.style={only marks, mark=triangle*, mark size=3.2pt, draw=iyiRenk, fill=iyiRenk,
    /pgfplots/error bars/x dir=both, /pgfplots/error bars/x explicit, /pgfplots/error bars/error bar style={line width=1.4pt, iyiRenk}, /pgfplots/error bars/error mark options={rotate=90, mark size=2.6pt, iyiRenk}},
  fseven last/.style={only marks, mark=square, mark size=2.3pt, draw=black!60, fill=white, line width=0.9pt,
    /pgfplots/error bars/x dir=both, /pgfplots/error bars/x explicit, /pgfplots/error bars/error bar style={line width=1.0pt, black!60, dashed}, /pgfplots/error bars/error mark options={rotate=90, mark size=2.6pt, black!60}},
}
\begin{tikzpicture}
% ---------------- panel (a): YOLO26n
\begin{axis}[name=A, fseven, height=6.4cm, yticklabels from table={\figroot data/f7a_yolo26n.dat}{label},
  title={(a) YOLO26n, KITTI test (13 scene groups)}, xticklabels={}]
\draw[line width=1.1pt, black] (axis cs:0,0.4) -- (axis cs:0,7.6);
\addplot[draw=none, mark=none] table[x expr=0, y=ypos]{\figroot data/f7a_yolo26n.dat};% ordered full y list for ytick=data
\addplot[fseven best, y filter/.expression={\thisrow{isgood}==0 ? y : nan}] table[x=best_m, y expr=\thisrow{ypos}+0.17, x error=best_sd]{\figroot data/f7a_yolo26n.dat};
\addplot[fseven good, y filter/.expression={\thisrow{isgood}==1 ? y : nan}] table[x=best_m, y expr=\thisrow{ypos}+0.17, x error=best_sd]{\figroot data/f7a_yolo26n.dat};
\addplot[fseven last] table[x=last_m, y expr=\thisrow{ypos}-0.17, x error=last_sd]{\figroot data/f7a_yolo26n.dat};
\end{axis}
% ---------------- panel (b): YOLO11s
\begin{axis}[name=B, fseven, at={(A.south west)}, anchor=north west, yshift=-1.4cm, height=6.4cm, yticklabels from table={\figroot data/f7b_yolo11s.dat}{label},
  title={(b) YOLO11s, KITTI test (13 scene groups)}, xlabel={$\Delta$ test mAP50-95 vs.\ the same-seed baseline [points]},
  legend style={at={(0.5,-0.24)}, anchor=north, legend columns=2, font=\footnotesize, draw=none, /tikz/every even column/.append style={column sep=8pt}}]
\draw[line width=1.1pt, black] (axis cs:0,0.4) -- (axis cs:0,7.6);
\addplot[draw=none, mark=none] table[x expr=0, y=ypos]{\figroot data/f7b_yolo11s.dat};% ordered full y list for ytick=data
\addplot[fseven best, y filter/.expression={\thisrow{isgood}==0 ? y : nan}] table[x=best_m, y expr=\thisrow{ypos}+0.17, x error=best_sd]{\figroot data/f7b_yolo11s.dat};
\addplot[fseven good, y filter/.expression={\thisrow{isgood}==1 ? y : nan}] table[x=best_m, y expr=\thisrow{ypos}+0.17, x error=best_sd]{\figroot data/f7b_yolo11s.dat};
\addplot[fseven last] table[x=last_m, y expr=\thisrow{ypos}-0.17, x error=last_sd]{\figroot data/f7b_yolo11s.dat};
\addlegendimage{only marks, mark=*, mark size=2.4pt, black}\addlegendentry{\texttt{best.pt}, mean $\pm$1 s.d. (3 seeds)}
\addlegendimage{only marks, mark=square, mark size=2.3pt, black!60, fill=white, line width=0.9pt}\addlegendentry{\texttt{last.pt}, mean $\pm$1 s.d. (3 seeds)}
\addlegendimage{only marks, mark=triangle*, mark size=3.2pt, iyiRenk}\addlegendentry{better, Holm-corrected $p<0.10$, $\Delta\ge+0.40$}
\end{axis}
\end{tikzpicture}
\caption{KITTI test: per-card $\Delta$ mAP50-95 vs.\ baseline [points]. Points show the three-seed mean $\pm$1 s.d. for \texttt{best.pt} and \texttt{last.pt}. (a) YOLO26n, (b) YOLO11s. Blue triangle: Holm-5 acceptance.}\label{fig:best-last}
\end{figure}
}

\paragraph{\texttt{best.pt} and \texttt{last.pt} per card.}
Fig.~\ref{fig:best-last} plots the per-card \texttt{best.pt} and \texttt{last.pt} $\Delta$s side by side. For YOLO11s (b), the \texttt{best.pt} $\Delta$ is positive in seven of seven cards, including H02 and H05, which are downward on VOC. The \texttt{last.pt} $\Delta$s, by contrast, scatter around zero. For YOLO26n (a), the pattern is mirrored. Three cards are negative at \texttt{best.pt}, whereas all seven are above the baseline at \texttt{last.pt}. The robustness label was ``sensitive to \texttt{best.pt} selection'' for 7 of the 14 cards (Table~\ref{tab:kitti-test} of the main article). For YOLO11s, H02 and H05 are labeled ``selection-independent''. This is a borderline case of the sign rule, because their \texttt{last.pt} $\Delta$s are $\approx$0. The common shift of the cards is the mean test mAP50-95 of the seven cards minus the baseline mean. For YOLO11s, it is +1.38 points at \texttt{best.pt} and $-$0.27 at \texttt{last.pt}. For YOLO26n, it is +0.01 and +0.65. The shift thus depends on the baseline checkpoint. The seven cards are measured against a single baseline, so 7/7 is not seven independent pieces of evidence. For YOLO11s, the clustered CI excludes zero for five cards. This too fits the common shift better than a card effect. Remaining hypothesis: the validation set is small (10 scene groups). On it, \texttt{best.pt} selection captures a peak specific to this set, and for YOLO11s it often picks the epoch-1 weights. This peak partly does not carry over to the test set (winner's curse and scene-group noise). The other explanations and the training pipeline diagnosis are given in the paragraphs below. The experiment that would discriminate among the explanations was not run. It would require fixed-epoch evaluation at intermediate epochs, or more seeds.

\paragraph{Epoch-1 weights.}
For YOLO11s, the best epoch is 1 in 11 of 24 runs and below 20 in 16 runs. The median is 10.5 (Table~\ref{tab:kitti-epoch}). For YOLO26n, no run has a best epoch of 1 (0/24), 10 runs are below 20, and the median is 24. In YOLO11s seed 2, \texttt{best.pt} is the epoch-1 weights in six of the eight arms: baseline, H01, H02, H03, H04, and H06. Five of the card--baseline comparisons are therefore between two epoch-1 weights. Their mean test $\Delta$ is +1.235, with a range of +0.634 to +1.593. These five comparisons are not independent, because they share the same baseline weights. The test mAP50-95 of baseline seed 2 is 18.22 at \texttt{best.pt} and 19.83 at \texttt{last.pt}.

\paragraph{Training pipeline.}
We measured mAP50-95 on 1000 training images. For YOLO26n baseline seed 0, \texttt{best.pt} gives 56.06 and \texttt{last.pt} gives 78.42. The validation value of the same \texttt{best.pt} is 17.89. For YOLO11s baseline seeds 0--2, \texttt{last.pt} gives 85.13, 85.00, and 85.25. The validation value of the same weights is 14.8--15.9. The epoch-1 \texttt{best.pt} has not yet learned the training scenes: it gives 36.74 for YOLO11s seed 2. Seeds 0 and 1, whose best epochs are 9 and 32, give 49.84 and 69.93. The low validation and test levels are consistent with a generalization gap rather than a pipeline error. One possible cause is overfitting to the 36 training scene groups. Another is that the optimizer and learning rate (MuSGD, 0.01) do not suit these data; they were fixed to preserve the VOC behavior. A third possibility is that full fine-tuning degrades the pretrained features on the scene-disjoint validation and test sets. This effect has been shown under large distribution shift \cite{kumar_2022_finetuning}. However, the magnitude of the shift between the KITTI splits was not measured. None of the three possibilities was tested.

\paragraph{Selection criterion mismatch.}
The mismatch hypothesis assumes that detections in \texttt{DontCare} regions penalize late epochs. It predicts that the best$-$last difference narrows under the \texttt{DontCare} rule. We measured the same weights on validation with this study's COCO-style evaluator, both without \texttt{DontCare} and with the \texttt{DontCare} rule. The difference widens in 48 of 48 runs. The mean goes from +3.23 to +3.62 for YOLO26n and from +3.84 to +4.31 for YOLO11s. The hypothesis is not supported. On validation, \texttt{best.pt} beats \texttt{last.pt} in 48/48 runs. This is structural, because \texttt{best.pt} is selected on validation, and it does not discriminate between explanations. Ultralytics' in-training evaluator (without \texttt{DontCare}) is a separate measurement. \texttt{best.pt} is selected by Ultralytics fitness, not by the evaluator of the primary endpoint (numbers in Table~\ref{tab:kitti-epoch}). On the test set, the baseline \texttt{best.pt} $-$ \texttt{last.pt} difference is +1.50 for YOLO26n and $-$0.36 for YOLO11s. On validation with the \texttt{DontCare} rule, the same difference is +3.78 and +3.43 points. For YOLO11s, the validation peak has no counterpart on the test set. For YOLO26n, it partly carries over.

\paragraph{Common shift.}
The common shift of the cards relative to the baseline is the mean test mAP50-95 of the seven cards minus the baseline mean. For YOLO11s, it is +1.38 points at \texttt{best.pt} and $-$0.27 at \texttt{last.pt}, a difference of 1.65. For YOLO26n, the pattern is mirrored: the shift is +0.01 at \texttt{best.pt} and +0.65 at \texttt{last.pt} (Fig.~\ref{fig:best-last}a). The common shift depends on which checkpoint represents the baseline. The seven cards are measured against a single baseline, so 7/7 is not seven independent pieces of evidence. The outlier in seed 2 does not explain the pattern on its own. For YOLO11s, seven of seven cards are also positive in seeds 0 and 1. The per-seed common shift is +1.36 and +0.78 there, and +1.99 in seed 2. The sign rule labels H02 and H05 on YOLO11s as ``selection-independent''. However, the \texttt{last.pt} $\Delta$s of these two cards are zero (+0.00, +0.03). Their magnitude at \texttt{best.pt} (+1.14, +2.28) depends on the selection. For these two cards, the label is a borderline case of the sign rule.

\paragraph{Conclusion.}
The H01 acceptance is reported under the pre-specified rule. It carries the ``sensitive to \texttt{best.pt} selection'' label and a note that most of the gain comes from a single test scene group. No claim of a robust improvement is made. Measured: for YOLO11s, the \texttt{best.pt} and \texttt{last.pt} $\Delta$s diverge systematically. For YOLO26n, they diverge in the opposite direction. Remaining hypothesis: \texttt{best.pt} selection captures a peak specific to a small validation set of 10 scene groups. This peak partly does not carry over to the test set. The causes would be selection bias (``winner's curse'') and scene-group noise. The experiment that would separate these explanations was not run. It would require fixed-epoch evaluation at intermediate epochs, or more seeds.

\FloatBarrier
\subsection{KITTI robustness analysis (post-unlock, exploratory)}\label{sec:ek-saglamlik}
These analyses were performed after the KITTI test set was opened. They involve no new test opening, inference, or training. The existing \texttt{best.pt} and \texttt{last.pt} test predictions were re-read, and the $\Delta$s recorded on the full sample were reproduced. These analyses are not part of the pre-specified rules, and no Holm correction was applied.

% generator: evidence/scripts/build_ek_tables.py <- evidence/evidence.json (ek5.sag.*; manual edits are lost at the next regeneration)
\begin{table}[p]
\centering\small\setlength{\tabcolsep}{3pt}
\caption{KITTI test robustness (post-unlock, exploratory): card $-$ baseline $\Delta$ [points]. Hierarchical bootstrap over seeds and the 13 scene groups ($B=2000$). Blue+bold: the CI excludes zero.}\label{tab:kitti-saglamlik}
\begin{tabularx}{\textwidth}{>{\raggedright\arraybackslash}X c c c c c c}
\toprule
Card & $\Delta$ & Hierarchical CI95 & $P(\Delta\le0)$ & $\Delta$ w/o tram & 6-class $\Delta$ & Pos. groups \\
\midrule
\multicolumn{7}{@{}l}{\textbf{YOLO26n}} \\
H01 multi-scale & +0.12 & [$-$2.26, +2.51] & 0.401 & +0.36 & +0.11 & 5/13 \\
H02 Inner-CIoU & $-$0.93 & [$-$2.82, +0.85] & 0.778 & $-$0.28 & $-$1.21 & 5/13 \\
H03 box copy-paste & +1.01 & [$-$0.81, +3.36] & 0.093 & +1.78 & +1.20 & 7/13 \\
H04 mosaic scale & $-$1.24 & [$-$3.05, +1.65] & 0.758 & $-$0.38 & $-$1.81 & 9/13 \\
H05 Varifocal & +0.61 & [$-$1.53, +2.32] & 0.278 & +0.25 & +0.77 & 7/13 \\
H06 box loss gain & $-$0.25 & [$-$2.06, +1.54] & 0.600 & +0.09 & $-$0.31 & 5/13 \\
H08 high dose & +0.71 & [$-$0.93, +2.73] & 0.147 & +1.24 & +0.87 & 7/13 \\
\midrule
\multicolumn{7}{@{}l}{\textbf{YOLO11s}} \\
H01 multi-scale & +1.49 & [$-$0.67, +2.61] & 0.077 & +0.53 & +2.03 & 9/13 \\
H02 Inner-CIoU & +1.14 & [$-$0.76, +2.37] & 0.120 & +0.19 & +1.53 & 10/13 \\
H03 box copy-paste & +1.41 & [$-$0.36, +3.50] & 0.050 & +1.20 & +1.70 & 11/13 \\
H04 mosaic scale & +1.07 & [$-$1.00, +3.01] & 0.116 & +0.72 & +0.96 & 11/13 \\
H05 Varifocal & \iyi{+2.28} & \iyi{[+0.20, +4.10]} & \iyi{0.018} & +1.42 & +3.15 & 9/13 \\
H06 box loss gain & +0.57 & [$-$0.81, +2.68] & 0.150 & +0.82 & +0.68 & 10/13 \\
H08 high dose & +1.68 & [$-$0.45, +4.62] & 0.084 & +1.54 & +1.78 & 11/13 \\
\bottomrule
\addlinespace[3pt]
\multicolumn{7}{@{}p{\dimexpr\textwidth-2\tabcolsep\relax}@{}}{$P(\Delta\le0)$: share of bootstrap draws with $\Delta\le0$. W/o tram: 7 classes. 6 classes: excluding person\_sitting and misc. These two columns are uncolored (no CI is shown). Pos. groups: scene groups with seed-mean $\Delta>0$. All values use \texttt{best.pt}. H02, H05: directional cards (expected $\Delta<0$).} \\
\end{tabularx}
\end{table}


Table~\ref{tab:kitti-saglamlik} reports a hierarchical bootstrap. Each replicate resamples with replacement, first the three seeds and then the 13 test scene groups. The group draw is shared by all cards and seeds (paired). The seed draw is over only $n=3$, so the intervals may still be optimistic. The interval covers zero for 13 of the 14 cards. The YOLO11s H01 numbers are given in Section~\ref{sec:genelleme-sonuc} of the main article. Only one card has an eight-class $\Delta$ interval that excludes zero (blue): YOLO11s H05, with [+0.20, +4.10]. However, H02 and H05 are directional replication cards, with an expected $\Delta<0$. In addition, the \texttt{last.pt} $\Delta$ of H05 is $\approx$0. The blue cell is therefore not an improvement claim. The without-tram and 6-class columns show no interval, so they are not colored. For the without-tram $\Delta$, the interval is [+0.11, +3.32] for YOLO11s H05 and [+0.07, +3.34] for H08. The positive-group column is not image-weighted. Single-group $\Delta$s are unstable for small groups of 5--6 images.

The YOLO11s H01 acceptance has further sensitivities. The one-sided $p$ of the clustered bootstrap is 0.0535. This is a sensitivity analysis, not a pre-specified test. The acceptance relies on the normality assumption of the $t$ test. With three seeds, a sign-flip test cannot give a one-sided $p$ below 0.125. The rule document does not provide for pooling across models. As a sensitivity analysis, we corrected the 10 improvement tests of the two models jointly with Holm-10, which gives $p=0.0848$. Pooling all 14 tests into a single family gives $p=0.119$. These values do not change the pre-specified decision. Leaving out drive\_0056 lowers the $\Delta$ of H01 from +1.49 to +0.45, a drop of $\approx$70\%. mAP is not additive over groups, so this is a leave-out difference. The gain does not appear in the standard three classes (car, pedestrian, cyclist). It comes from the classes with few samples (Table~\ref{tab:kitti-ikincil}).

The clustered bootstrap CI measures only the sampling uncertainty of the test scene groups at fixed weights. It does not include seed variance. For the improvement cards that were not accepted, its upper bound is at most +2.92 for YOLO26n and +2.76 for YOLO11s. For YOLO11s, it excludes zero for five of seven cards (H03, H04, H05, H06, H08). For the accepted H01, in contrast, it crosses zero. For YOLO26n, it excludes zero for no card. This pattern points to a common difference that shifts all cards against the same baseline, not to a card-specific effect (Section~\ref{sec:ek-kontrol-noktasi}). It is not taken as evidence of a card effect.

Validation (987 images, 10 scene groups) is a secondary endpoint. It did not select the same card as the test set. For YOLO26n, there is no acceptance on validation either. For YOLO11s, only H04 met the acceptance criterion on validation, with +1.82 $\pm$ 0.43 and $p_{\mathrm{Holm}}=0.046$. This is not a decision. H01 was not accepted on validation, with +1.73 $\pm$ 1.78 and $p_{\mathrm{Holm}}=0.350$. On the test set, the picture reverses. H01 meets the criterion, whereas H04 does not ($+1.07\pm0.72$, $p_{\mathrm{Holm}}=0.252$). The two sets select different cards. This is consistent with the identity of the single accepted YOLO11s card being sensitive to scene-group sampling.

\subsection{Details of training-free inference (H07)}\label{sec:ek-kitti-h07}
H07 was run on validation with the baseline \texttt{best.pt} weights. For each model, the single row with the highest validation mean entered the test. For YOLO11s seed 2, these \texttt{best.pt} weights are the epoch-1 weights (Section~\ref{sec:ek-kontrol-noktasi}).

For YOLO26n, the selected U832 row gave +1.50 on validation and +1.54 $\pm$ 0.58 points on the test set. Latency at batch 1 rose from 7.74 ms to 7.80 ms, a ratio of 1.008. This is within the pre-specified $\times1.10$ budget. The pixel count grew 1.69-fold, yet latency increased by only 0.8\%. This suggests that per-image Python and preprocessing overhead dominate the measurement. On VOC, the latency ratio of the same row was 1.14.

For YOLO11s, the selected W2 row is the WBF fusion of the T2 branches. It gave +1.61 on validation and +2.02 $\pm$ 0.95 on the test set. Latency rose from 6.82 ms to 26.70 ms ($\times3.91$), which is outside the pre-specified $\times1.10$ budget.

The two models use different latency engines. YOLO11s uses the VOC engine, and YOLO26n uses per-image \texttt{predict}. At R0, nano gives 7.74 ms and small gives 6.82 ms. Absolute ms values are therefore not compared across models; only the within-model ratio is read. These $\Delta$s are measured with the same weights. They therefore do not include the \texttt{best.pt} selection difference between the cards and the baseline. However, they are single-row, three-seed, and exploratory, and no CI or Holm correction is given. On VOC validation, 832 px input had changed mAP50-95 by $-$5.19 (Section~\ref{sec:bulgu-gecikme} of the main article). On KITTI, the same row has a positive sign. The untested hypothesis for this contrast is the resolution difference described above (Section~\ref{sec:ek-kitti-kurulum}).

\paragraph{Latency measurement error and its correction.}
The first run of the KITTI H07 script crashed with a CUDA out-of-memory error. This happened during the latency measurement of the selected row, and no latency value was written. The image paths had been passed to \texttt{predict} as a single source list. Ultralytics processes such a list in a single batch, so the batch size of 1 required by the rule document was not met. The error was corrected before any latency value was produced and before the KITTI test set was opened. For YOLO11s, the VOC H07 measurement engine was used unchanged. For YOLO26n, per-image \texttt{predict} with GPU synchronization was used. The 50 warm-up iterations, 300 measurements, and FP16 were unchanged. The definition of the accuracy rows and the selection rule were also unchanged. The reported KITTI latencies come from the corrected measurement. The VOC latencies were not affected.

\FloatBarrier
